\section{\texorpdfstring{$C^\infty(M)$}{C-infinity(M)} 上的不变范数与霍费尔长度}
\stfhIndex{I170}{bu^bian^fan^shu}{不变范数（invariant norm）}
\label{sec:invariant-norms-hofer-length}

$C^\infty(M)/\mathbb R$ 上有一个自然的 $\Diff(M)$-不变范数，定义为
\begin{equation}
 \|h\|=\operatorname{osc}(h)=\max h-\min h.
 \label{eq:oscillation-norm}
\end{equation}
事实表明，在下面的意义下，当 $n=\dim M\geq2$ 时，范数
\eqref{eq:oscillation-norm} 是\emph{唯一的} $\Diff(M)$-不变范数 (Che00)。

\begin{theorem}[切卡诺夫]\label{thm:chekanov-invariant-norm}
\stfhIndex{I074}{qie^ka^nuo^fu}{切卡诺夫（Chekanov）}
设 $\dim M\geq2$。假设 $\|\mathord\cdot\|'$ 是
$C_0^\infty(M)$ 上任意 $\Diff(M)$-不变范数。则存在 $\lambda\geq0$，使得对所有 $f\in C_0^\infty(M)$，
\[
 \|f\|'=\lambda\|f\|.
\]
\end{theorem}

当考虑辛流形 $(M,\omega)$ 时，可以根据 $M$ 是闭的还是开的，采用上面提到的归一化截面，把商空间 $C^\infty(M)/\mathbb R$ 实现为
$C^\infty(M)$ 的一个不变子集。令 $G$ 是两个群
$\Ham^{sc}(M,\omega)$ 与 $\Ham^{wc}(M,\omega)$ 中的任意一个。注意，当 $M$ 闭时，因为 $\partial M=\varnothing$，有
$\Ham(M,\omega)=\Ham^{sc}(M,\omega)$。然而，当 $M$ 闭时，我们取
$C_0^\infty(M)$ 作为相应的李代数；当 $M$ 开时，则取
$C_c^\infty(M)$。

上面切卡诺夫定理的如下辛（或哈密顿）对应版本，是一个重要的问题。

\begin{question}\label{q:classification-ham-invariant-norms}
设 $(M,\omega)$。相应李代数 $\mathfrak g$ 上的 $G$-不变范数如何分类？
\end{question}

亚龙·奥斯特罗弗（Yaron Ostrover）与罗伊·瓦格纳（Roy Wagner） (OW05) 以及布霍夫斯基与奥斯特罗弗 (BO13) 所得结果的结合，给出了这个问题一个相当完整的答案。

\begin{theorem}[奥斯特罗弗--瓦格纳]\label{thm:ostrover-wagner-pseudodistance}
假设 $(M,\omega)$ 是闭的。设 $|\mathord\cdot|$ 是
$C_0^\infty(M)$ 上的 $\Ham(M,\omega)$-不变范数，并且对某个常数
$C$ 有 $|\mathord\cdot|\leq C|\mathord\cdot|_\infty$，但这两个范数不等价。则 $\Ham(M,\omega)$ 上相伴的芬斯勒伪距离恒等于零。
\end{theorem}

\begin{theorem}[布霍夫斯基--奥斯特罗弗]\label{thm:buhovsky-ostrover-pseudonorm}
对闭辛流形 $(M,\omega)$，$C_0^\infty(M)$ 上任意
$\Ham(M,\omega)$-不变伪范数 $|\mathord\cdot|$ 都被 $L_\infty$ 范数从上控制；即对某个常数 $C\geq0$，有
$|\mathord\cdot|\leq C|\mathord\cdot|_\infty$。
\end{theorem}

于是，上述不变范数以下面的方式在李群 $G$ 上诱导一个芬斯勒型度量。对 $\phi\in G$，可以借助赋值
\[
 f\longmapsto f\circ\phi,
 \qquad f\in C_0^\infty(M),
\]
把 $G$ 在 $\phi$ 处的切空间与 $\mathfrak g$ 本身等同。把 $f$ 看作
李代数 $\mathfrak g$ 中的元素，我们可将赋值
\[
 \phi\longmapsto f\circ\phi
\]
看作 $f$ 生成的 $G$ 上的右不变向量场。

为简化叙述，我们把讨论限制在 $M$ 闭的情形；稍后将简要指出开情形的差别。

设 $\lambda:[0,1]\to\Symp(M,\omega)$ 是一条始于单位元的哈密顿路径。回忆定义 2.3.15 中的记号
\stfhIndex{I014}{00-symbol-014}{$\operatorname{Dev}(\lambda)$}
\stfhIndex{I045}{00-symbol-045}{$\operatorname{Tan}(\lambda)$}
\[
 \operatorname{Tan}(\lambda)(t,x)
 :=H(t,\phi_H^t(x)),
 \qquad
 \operatorname{Dev}(\lambda)(t,x):=H(t,x).
\]
现在按标准方式，关于上述芬斯勒结构定义 $\Ham(M,\omega)$ 上路径
$\lambda$ 的芬斯勒长度。

\begin{definition}\label{def:hofer-length}
\stfhIndex{I027}{00-symbol-027}{$\operatorname{leng}(\lambda)$}
设 $\lambda\in\mathcal P^{\mathrm{ham}}(\Symp(M,\omega))$，并设 $H$ 是生成给定哈密顿路径 $\lambda$ 的哈密顿量，即
$\lambda(t)=\phi_H^t\circ\lambda(0)$。哈密顿路径 $\lambda$ 的\emph{霍费尔长度}定义为
\begin{equation}
 \operatorname{leng}(\phi_H)
 =\int_0^1\left(
   \max_x H(t,\phi_H^t(x))-
   \min_x H(t,\phi_H^t(x))\right)dt.
 \label{eq:hofer-length}
\end{equation}
\end{definition}

显然，$\operatorname{leng}(\phi_H)$ 与熟知的 $L^{(1,\infty)}$ 范数表达式一致：
\[
 \|H\|=\int_0^1\operatorname{osc}H_t\,dt
 =\int_0^1\left(\max_xH_t-\min_xH_t\right)dt.
\]
下面几个简单练习在稍后计算霍费尔长度函数时会有用。

\begin{exercise}\label{ex:hofer-length-identities}
设 $H,K:[0,1]\times M\to\mathbb R$。证明：
\begin{enumerate}[label=(\arabic*)]
 \item $\operatorname{leng}((\phi_H^1)^{-1}\phi_K)
 =\operatorname{leng}(\phi_K^{-1}\phi_H)$；
 \item $\operatorname{leng}(\phi_H\phi_K)
 \leq\operatorname{leng}(\phi_H)+\operatorname{leng}(\phi_K)$；
 \item $\operatorname{leng}(\phi_H)
 =\operatorname{leng}((\phi_H^1)^{-1})$。
\end{enumerate}
（提示：这些证明本质上是命题 2.3.8 的平移。）
\end{exercise}

长度函数按如下方式在
$\mathcal P^{\mathrm{ham}}(\Symp(M,\omega),\operatorname{id})$ 上给出一个自然拓扑。

\begin{definition}\label{def:hofer-topology-path-space}
\stfhIndex{I039}{00-symbol-039}{$\mathcal P^{\mathrm{ham}}(\Symp(M,\omega),\operatorname{id})$}
\stfhIndex{I162}{huo^fei^er^tuo^pu}{霍费尔拓扑（Hofer topology）}
设
$\lambda,\mu\in\mathcal P^{\mathrm{ham}}(\Symp(M,\omega),\operatorname{id})$。定义
$\mathcal P^{\mathrm{ham}}(\Symp(M,\omega),\operatorname{id})$ 上的距离函数为
\begin{equation}
 d(\lambda,\mu):=\operatorname{leng}(\lambda^{-1}\circ\mu),
 \qquad
 \lambda,\mu\in
 \mathcal P^{\mathrm{ham}}(\Symp(M,\omega),\operatorname{id}),
 \label{eq:hofer-path-distance}
\end{equation}
其中 $\lambda^{-1}\circ\mu$ 是哈密顿路径
$t\in[0,1]\mapsto\lambda(t)^{-1}\mu(t)$。称诱导的拓扑为
$\mathcal P^{\mathrm{ham}}(\Symp(M,\omega),\operatorname{id})$ 的\emph{霍费尔拓扑}。
\end{definition}
